\documentclass[a4paper]{article}

% ─── Encodage & langue ───────────────────────────────────────────────────────
\usepackage[utf8]{inputenc}
\usepackage[T1]{fontenc}
\usepackage[french]{babel}

% ─── Mise en page ────────────────────────────────────────────────────────────
\usepackage[top=3cm, bottom=3cm, left=3cm, right=3cm]{geometry}
\setlength{\parindent}{1.5em}
\setlength{\parskip}{0.4em}

% ─── Typographie ─────────────────────────────────────────────────────────────
\usepackage{microtype}

% ─── Couleurs ────────────────────────────────────────────────────────────────
\usepackage{xcolor}

% Palette info (bleu doux)
\definecolor{InfoBack}{RGB}{235, 244, 255}
\definecolor{InfoRule}{RGB}{59,  130, 246}
% Palette attention (ambre)
\definecolor{WarnBack}{RGB}{255, 249, 219}
\definecolor{WarnRule}{RGB}{217, 149,  0}
% Palette étape (gris neutre)
\definecolor{StepBack}{RGB}{246, 247, 248}
\definecolor{StepHead}{RGB}{74,  85,  104}
% Code
\definecolor{CodeBack}{RGB}{248, 249, 250}
\definecolor{CodeRule}{RGB}{200, 205, 210}
% Sections
\definecolor{SecColor}{RGB}{30,  64, 120}

% ─── Boîtes ──────────────────────────────────────────────────────────────────
\usepackage{tcolorbox}
\tcbuselibrary{skins, breakable}

% Boîte « note / info » — fond bleu très clair, règle gauche bleue
\newtcolorbox{infobox}{
  enhanced, breakable,
  colback=InfoBack,
  colframe=InfoRule,
  leftrule=3.5pt, rightrule=0.3pt, toprule=0.3pt, bottomrule=0.3pt,
  arc=3pt,
  left=8pt, right=8pt, top=5pt, bottom=5pt,
  fontupper=\small,
}

% Boîte « attention / warning » — fond ambre très clair, règle gauche ambre
\newtcolorbox{warnbox}{
  enhanced, breakable,
  colback=WarnBack,
  colframe=WarnRule,
  leftrule=3.5pt, rightrule=0.3pt, toprule=0.3pt, bottomrule=0.3pt,
  arc=3pt,
  left=8pt, right=8pt, top=5pt, bottom=5pt,
  fontupper=\small,
}

% Boîte « étape » — fond gris clair, bandeau de titre gris foncé
\newtcolorbox{stepbox}[1]{
  enhanced, breakable,
  colback=StepBack,
  colframe=StepHead,
  colbacktitle=StepHead,
  coltitle=white,
  title={\normalfont\small\bfseries #1},
  arc=3pt,
  leftrule=0.5pt, rightrule=0.5pt, toprule=0.5pt, bottomrule=0.5pt,
  left=8pt, right=8pt, top=5pt, bottom=5pt,
  attach boxed title to top left={xshift=10pt, yshift=-\tcboxedtitleheight/2},
  boxed title style={arc=2pt, boxrule=0pt},
  fontupper=\small,
}

% ─── Code source ─────────────────────────────────────────────────────────────
\usepackage{listings}

\lstdefinestyle{codestyle}{
  backgroundcolor=\color{CodeBack},
  basicstyle=\ttfamily\small,
  commentstyle=\color{gray}\itshape,
  breaklines=true,
  keepspaces=true,
  numbers=left,
  numberstyle=\tiny\color{gray},
  numbersep=10pt,
  showstringspaces=false,
  tabsize=2,
  frame=single,
  rulecolor=\color{CodeRule},
  framesep=5pt,
  xleftmargin=2em,
  framexleftmargin=2em,
  aboveskip=8pt,
  belowskip=8pt,
}

\lstnewenvironment{bashcode}{\lstset{style=codestyle}}{}
\lstnewenvironment{yamlcode}{\lstset{style=codestyle}}{}

% ─── Liens hypertexte ────────────────────────────────────────────────────────
\usepackage{hyperref}
\hypersetup{
  colorlinks=true,
  linkcolor=SecColor,
  urlcolor=SecColor,
  pdftitle={Projet Cloud Étudiant},
  pdfauthor={Infrastructure Cloud},
}

% ─── En-tête / pied de page ──────────────────────────────────────────────────
\usepackage{fancyhdr}
\pagestyle{fancy}
\fancyhf{}
\fancyhead[L]{\small\color{gray}\itshape Projet Cloud Étudiant}
\fancyhead[R]{\small\color{gray}\itshape Infrastructure \& Orchestration}
\fancyfoot[C]{\small\color{gray}\thepage}
\renewcommand{\headrulewidth}{0.4pt}
\renewcommand{\footrulewidth}{0pt}

% ─── Commandes utilitaires ───────────────────────────────────────────────────
\newcommand{\cmd}[1]{\texttt{#1}}
\newcommand{\file}[1]{\texttt{\textit{#1}}}

% ─── Sections ────────────────────────────────────────────────────────────────
\usepackage{titlesec}
\titleformat{\section}
  {\large\bfseries\scshape\color{SecColor}}
  {}
  {0em}
  {\thesection\quad}
  [\color{SecColor}\vspace{0.3ex}\hrule\vspace{0.3ex}]
\titleformat{\subsection}
  {\normalsize\bfseries\color{SecColor!80!black}}
  {\thesubsection\quad}
  {0em}
  {}
\titleformat{\subsubsection}
  {\normalsize\itshape}
  {\thesubsubsection\quad}
  {0em}
  {}

% ═════════════════════════════════════════════════════════════════════════════
\begin{document}
% ═════════════════════════════════════════════════════════════════════════════

\begin{titlepage}
  \centering
  \vspace*{4cm}
  {\Large\bfseries\scshape Projet Cloud Étudiant\par}
  \vspace{0.5em}
  {\large Phase 1 --- Infrastructure et Orchestration\par}
  \vspace{3cm}
  \hrule
  \vspace{1cm}
  {\normalsize Infrastructure Cloud\par}
  \vspace{0.5em}
  {\small\itshape 2025\par}
  \vfill
\end{titlepage}

\tableofcontents
\newpage


% ═════════════════════════════════════════════════════════════════════════════
\section{Vue d'ensemble de l'infrastructure}
% ═════════════════════════════════════════════════════════════════════════════

Le projet Cloud Étudiant repose sur un empilement de trois couches technologiques qui coopèrent pour offrir une expérience de virtualisation en libre-service.

\begin{itemize}
  \item \textbf{Proxmox VE} (hyperviseur) : le serveur physique \cmd{exbasis} héberge les machines virtuelles et expose une API REST que les playbooks Ansible utilisent pour créer, configurer et supprimer les VMs.
  \item \textbf{AWX} (orchestrateur) : déployé dans une VM dédiée, AWX exécute les playbooks Ansible via une interface web et une API REST, rendant l'orchestration accessible sans manipuler la ligne de commande.
  \item \textbf{Apache Guacamole} (passerelle d'accès) : les étudiants accèdent à leurs VMs depuis un navigateur web, sans client SSH ni client RDP installé localement. Guacamole traduit les protocoles SSH (Linux) et RDP (Windows) en sessions web.
\end{itemize}

\subsection{Plan d'adressage réseau}

\begin{itemize}
  \item \textbf{Proxmox} : \cmd{172.16.1.61} (port d'administration : 8006)
  \item \textbf{Guacamole} : \cmd{172.16.180.250:8080/guacamole}
  \item \textbf{VMs étudiantes} : \cmd{172.16.205.X/16}, passerelle \cmd{172.16.1.213}\\
        Le dernier octet \cmd{X} est l'identifiant de l'étudiant (\cmd{vmid\_etudiant}), compris entre 10 et 99.
  \item \textbf{VMs Linux} : VMID Proxmox = \cmd{200X} (exemple : étudiant 12 \(\to\) VMID 20012)
  \item \textbf{VMs Windows} : VMID Proxmox = \cmd{300X} (exemple : étudiant 12 \(\to\) VMID 30012)
\end{itemize}

\begin{infobox}
  \textbf{Note.} Le nœud Proxmox se nomme \cmd{exbasis} dans tous les appels API. Ce nom figure explicitement dans les playbooks : toute migration vers un autre nœud nécessite de mettre à jour ce paramètre.
\end{infobox}

\subsection{Templates de VMs disponibles}

\begin{itemize}
  \item \textbf{UbuntuServ} \(\to\) template ID \cmd{106} (clone source : \cmd{VMtest})
  \item \textbf{Debian} \(\to\) template ID \cmd{103} (clone source : \cmd{VMtest})
  \item \textbf{Win10} \(\to\) template ID \cmd{107} (clone source : \cmd{WIN10-TEMPLATE})
\end{itemize}

\subsection{Flux d'un déploiement complet}

\begin{enumerate}
  \item L'étudiant (ou le portail web) déclenche le Job Template \file{deploy.yml} dans AWX en renseignant le Survey (identifiant, OS, services souhaités).
  \item AWX exécute \file{deploy.yml} dans un Execution Environment Kubernetes. Le playbook appelle l'API Proxmox via le module \cmd{community.proxmox.proxmox\_kvm} pour cloner la template.
  \item L'IP statique est injectée via l'API REST Proxmox, puis le disque Cloud-Init est régénéré.
  \item La VM démarre. AWX attend que SSH (Linux) ou RDP (Windows) soit accessible.
  \item Une connexion Guacamole est créée automatiquement, ainsi qu'un compte étudiant dédié.
  \item Le Job Template \file{services.yml} est lancé pour installer Apache et/ou MariaDB sur la VM.
\end{enumerate}


% ═════════════════════════════════════════════════════════════════════════════
\section{Configuration côté Proxmox}
% ═════════════════════════════════════════════════════════════════════════════

Avant qu'AWX puisse demander à Proxmox de cloner des machines virtuelles, il faut établir une relation de confiance entre les deux systèmes. Proxmox ne doit pas accepter n'importe quelle requête API : il convient de générer un \textbf{token d'API} dédié à AWX, qui agira comme une clé d'authentification spécifique.

\subsection{Pourquoi un token plutôt qu'un mot de passe ?}

On pourrait techniquement utiliser le couple login/mot de passe \cmd{root} de Proxmox, mais l'approche par token est nettement préférable pour trois raisons.

\begin{itemize}
  \item \textbf{Révocation ciblée.} Si le token est compromis, on le supprime sans avoir à changer le mot de passe \cmd{root}.
  \item \textbf{Traçabilité.} Les journaux Proxmox identifient les actions effectuées par token, ce qui permet de distinguer ce qu'AWX a fait de ce qu'un humain a fait directement.
  \item \textbf{Indépendance.} Changer le mot de passe \cmd{root} ne rompt pas l'automatisation.
\end{itemize}

\subsection{Création du token API dans Proxmox}

\begin{stepbox}{Étape 1 --- Accéder au gestionnaire de tokens}
  Se connecter à l'interface web de Proxmox (\cmd{https://172.16.1.61:8006}) avec le compte \cmd{root@pam}, puis naviguer dans le menu de gauche :

  \medskip
  \textbf{Datacenter \(\to\) Permissions \(\to\) API Tokens}

  \medskip
  Cliquer sur le bouton \textbf{Add} en haut du panneau central.
\end{stepbox}

\begin{stepbox}{Étape 2 --- Remplir le formulaire de création}
  \begin{itemize}
    \item \textbf{User} : \cmd{root@pam}
    \item \textbf{Token ID} : \cmd{awx} (identifiant libre, sans espace ni caractère spécial)
    \item \textbf{Expire} : laisser vide pour un token sans expiration
    \item \textbf{Privilege Separation} : \textbf{décocher cette case} --- lorsqu'elle est cochée (par défaut), Proxmox crée un token sans aucune permission, même si l'utilisateur est \cmd{root}. En la décochant, le token hérite directement des droits \cmd{root}, ce qui simplifie la mise en route dans un contexte pédagogique.
    \item \textbf{Comment} : \textit{Token utilisé par AWX pour l'orchestration des VMs étudiantes} (optionnel)
  \end{itemize}

  Cliquer sur \textbf{Add}.
\end{stepbox}

\begin{stepbox}{Étape 3 --- Sauvegarder le secret}
  Une fenêtre de confirmation affiche deux valeurs critiques :
  \begin{itemize}
    \item \textbf{Token ID complet} : \cmd{root@pam!awx}
    \item \textbf{Secret} : une chaîne au format UUID, par exemple \cmd{12345678-abcd-1234-abcd-1234567890ab}
  \end{itemize}
  \textbf{Copier immédiatement ces deux valeurs.} Le secret n'est affiché qu'une seule fois : si la fenêtre est fermée sans le copier, il faut supprimer le token et en créer un nouveau.
\end{stepbox}

\begin{warnbox}
  \textbf{Attention --- Sécurité du secret.} Le secret du token donne un accès équivalent à \cmd{root} sur Proxmox. Il ne doit jamais être commité dans Git, jamais écrit en dur dans un playbook. Dans AWX, il sera stocké chiffré dans un objet \textit{Credential}.
\end{warnbox}

\subsection{Utilisation du token dans les appels API directs}

Lorsque les playbooks appellent directement l'API REST Proxmox (pour injecter l'IP ou régénérer le disque Cloud-Init), le token est transmis dans l'en-tête HTTP \cmd{Authorization} :

\begin{bashcode}
Authorization: PVEAPIToken=root@pam!awx=<SECRET>
\end{bashcode}

Le module Ansible \cmd{community.proxmox.proxmox\_kvm} encapsule cela automatiquement à partir des paramètres \cmd{api\_user}, \cmd{api\_token\_id} et \cmd{api\_token\_secret} injectés par le credential AWX.


% ═════════════════════════════════════════════════════════════════════════════
\section{Déploiement de Guacamole}
% ═════════════════════════════════════════════════════════════════════════════

Apache Guacamole est la passerelle d'accès web qui permet aux étudiants de se connecter à leurs VMs depuis un navigateur, sans installer de client SSH ou RDP. Il est déployé en tant que conteneur LXC directement sur le serveur Proxmox, à l'adresse \cmd{172.16.180.250}.

La méthode retenue utilise les \textbf{Proxmox VE Helper Scripts} : une collection de scripts communautaires maintenus sur GitHub qui automatisent la création et la configuration de conteneurs LXC courants sur Proxmox. Pour Guacamole, le script crée un conteneur Debian, installe Tomcat, le serveur Guacamole, \cmd{guacd} (le démon de traduction de protocoles) et MariaDB en une seule commande.

\begin{infobox}
  \textbf{Note.} Les scripts sont hébergés sur \cmd{https://community-scripts.github.io/ProxmoxVE/}. Vérifier l'URL officielle avant d'exécuter un script téléchargé depuis Internet et en examiner le contenu au préalable.
\end{infobox}

\subsection{Création du conteneur LXC}

Depuis le \textbf{shell du nœud Proxmox} (\textit{Datacenter \(\to\) exbasis \(\to\) Shell}), exécuter :

\begin{bashcode}
bash -c "$(wget -qLO - \
  https://github.com/community-scripts/ProxmoxVE/raw/main/ct/guacamole.sh)"
\end{bashcode}

Le script affiche un assistant interactif. Les choix recommandés pour ce projet :

\begin{itemize}
  \item \textbf{Default Settings} ou \textbf{Advanced} : choisir \textit{Advanced} pour fixer l'IP statique.
  \item \textbf{IP address} : \cmd{172.16.180.250/16} (adresse référencée dans tous les playbooks).
  \item \textbf{Gateway} : \cmd{172.16.1.213}
  \item \textbf{RAM} : 1024 Mo minimum, 2048 Mo recommandé.
  \item \textbf{Disk size} : 4 Go suffisent.
  \item \textbf{CPU cores} : 2.
\end{itemize}

Le script crée le conteneur, installe tous les composants et démarre les services automatiquement. L'opération prend environ deux à trois minutes.

\subsection{Accès initial et changement du mot de passe}

Une fois le conteneur démarré, Guacamole est accessible via :

\begin{bashcode}
http://172.16.180.250:8080/guacamole
\end{bashcode}

Les identifiants par défaut sont \cmd{guacadmin} / \cmd{guacadmin}. \textbf{Changer immédiatement ce mot de passe} :

\begin{enumerate}
  \item Se connecter avec \cmd{guacadmin / guacadmin}.
  \item Cliquer sur \textit{guacadmin} en haut à droite \(\to\) \textbf{Settings \(\to\) Preferences}.
  \item Saisir le mot de passe actuel et définir un nouveau mot de passe.
  \item Cliquer sur \textbf{Update Password}.
\end{enumerate}

\begin{warnbox}
  \textbf{Attention.} Les playbooks utilisent les identifiants \cmd{guacadmin / guacadmin} en dur dans les variables de \file{deploy.yml}. Si le mot de passe est changé, il faut mettre à jour les variables \cmd{guac\_admin} et \cmd{guac\_pass} dans le playbook, ou mieux, les passer via un credential AWX dédié.
\end{warnbox}

\subsection{Vérification des services dans le conteneur}

Pour s'assurer que tous les composants sont opérationnels, se connecter au shell du conteneur LXC depuis Proxmox (\textit{exbasis \(\to\) Guacamole LXC \(\to\) Console}) :

\begin{bashcode}
# Vérifier l'état des trois services principaux
systemctl status tomcat9
systemctl status guacd
systemctl status mariadb

# Vérifier que le port 8080 est en écoute
ss -tlnp | grep 8080
\end{bashcode}

Les trois services doivent être en état \cmd{active (running)}. Si \cmd{tomcat9} ne démarre pas, consulter les journaux :

\begin{bashcode}
journalctl -u tomcat9 -n 50
\end{bashcode}

\subsection{Rôle de Guacamole dans le projet}

Guacamole n'est pas géré manuellement après son installation initiale : c'est AWX qui administre ses ressources (connexions et comptes utilisateurs) via la collection Ansible \cmd{scicore.guacamole}. Concrètement :

\begin{itemize}
  \item À chaque déploiement de VM, \file{deploy.yml} appelle \cmd{guacamole\_connection} pour créer une entrée SSH ou RDP pointant vers la nouvelle VM.
  \item Un compte étudiant \cmd{etu\_\{\{ vmid\_etudiant \}\}} est créé (ou mis à jour) avec les droits d'accès à cette connexion.
  \item À la suppression de la VM, \file{suppression.yml} supprime la connexion correspondante pour maintenir Guacamole propre.
\end{itemize}

L'étudiant se connecte donc à Guacamole avec son compte personnel et n'a accès qu'à sa propre VM, sans voir celles des autres.


% ═════════════════════════════════════════════════════════════════════════════
\section{Déploiement de l'orchestrateur AWX}
% ═════════════════════════════════════════════════════════════════════════════

AWX est l'orchestrateur retenu pour ce projet. C'est la version libre de Red Hat Ansible Automation Platform, utilisée comme référence en entreprise pour piloter des playbooks Ansible à grande échelle.

AWX ne s'installe plus par Docker Compose depuis 2021. La méthode officielle utilise l'\textbf{AWX Operator} sur un cluster Kubernetes. Pour un déploiement sur un seul nœud, \textbf{K3s} est utilisé : distribution Kubernetes allégée, parfaitement adaptée à ce cas d'usage.

\begin{infobox}
  \textbf{Configuration matérielle recommandée pour la VM AWX.} 4 vCPU minimum, 8 Go de RAM, 40 Go de disque, Debian 12 à jour.
\end{infobox}

\subsection{Préparation du système}

\begin{bashcode}
apt update && apt upgrade -y
apt install git curl -y
\end{bashcode}

\subsection{Installation de Kubernetes (K3s)}

\begin{bashcode}
# Installation de K3s
curl -sfL https://get.k3s.io | sh -

# Passer en root
su -

# Déclarer le fichier de configuration Kubernetes
export KUBECONFIG=/etc/rancher/k3s/k3s.yaml

# Rendre cette configuration permanente
echo "export KUBECONFIG=/etc/rancher/k3s/k3s.yaml" >> ~/.bashrc

# Vérifier que le nœud est en état "Ready"
kubectl get nodes
\end{bashcode}

\subsection{Installation de l'AWX Operator}

L'Operator gère le cycle de vie complet d'AWX : déploiement, mise à jour, sauvegarde. La version \textbf{2.12.0} est utilisée pour éviter un bug de l'image \cmd{kube-rbac-proxy} présent sur les versions ultérieures.

\begin{bashcode}
# Créer l'espace de nommage isolé
kubectl create namespace awx

# Déployer l'opérateur AWX version 2.12.0
kubectl apply -k \
  "github.com/ansible/awx-operator/config/default?ref=2.12.0" \
  -n awx

# Vérifier (attendre que le pod soit "Running")
kubectl get pods -n awx
\end{bashcode}

\begin{warnbox}
  \textbf{Attention.} Si \cmd{kubectl get pods -n awx} n'affiche rien immédiatement, c'est normal : Kubernetes télécharge l'image. Patienter une à deux minutes et relancer la commande.
\end{warnbox}

\subsection{Déploiement de l'instance AWX}

\begin{bashcode}
# Créer le fichier de description de l'instance
cat > awx-instance.yaml << 'EOF'
apiVersion: awx.ansible.com/v1beta1
kind: AWX
metadata:
  name: awx-cloud
  namespace: awx
spec:
  service_type: nodeport
EOF

# Appliquer la configuration
kubectl apply -f awx-instance.yaml

# Suivre l'installation (5 à 10 minutes)
kubectl get pods -n awx -w
\end{bashcode}

L'option \cmd{service\_type: nodeport} expose AWX sur un port aléatoire (entre 30000 et 32767), accessible directement depuis l'extérieur sans avoir à configurer un Ingress.

\subsection{Récupération des accès}

Une fois les trois pods \cmd{awx-cloud-web}, \cmd{awx-cloud-task} et \cmd{awx-cloud-postgres-13-0} en état \cmd{Running} :

\begin{bashcode}
# Extraire et déchiffrer le mot de passe administrateur
kubectl get secret awx-cloud-admin-password \
  -o jsonpath="{.data.password}" -n awx | base64 --decode; echo

# Récupérer le port d'accès web (chercher la valeur 3XXXX)
kubectl get svc awx-cloud-service -n awx
\end{bashcode}

\begin{infobox}
  \textbf{Accès.} Ouvrir \cmd{http://<IP\_DE\_LA\_DEBIAN>:<PORT>} dans un navigateur. Identifiants par défaut : \cmd{admin} et le mot de passe récupéré ci-dessus.
\end{infobox}

\subsection{Les Execution Environments}

Point fondamental d'AWX : chaque playbook ne tourne pas directement sur la VM AWX, mais dans un \textbf{conteneur isolé} appelé \textit{Execution Environment} (EE). Ce conteneur embarque Ansible, Python et un ensemble de collections par défaut.

Conséquence directe : si un playbook a besoin d'une collection ou d'une bibliothèque Python supplémentaire, il faut le déclarer explicitement. Pour ce projet, le module \cmd{community.proxmox.proxmox\_kvm} dépend de la bibliothèque Python \cmd{proxmoxer}, absente de l'EE par défaut.

\begin{warnbox}
  \textbf{Attention.} L'EE de base n'inclut ni la collection \cmd{community.proxmox}, ni la bibliothèque \cmd{proxmoxer}. Deux solutions complémentaires sont mises en œuvre : déclarer les collections dans \file{collections/requirements.yml} (lu automatiquement par AWX avant chaque exécution), et installer \cmd{proxmoxer} via une tâche \cmd{pip} en tête de chaque playbook qui en a besoin.
\end{warnbox}


% ═════════════════════════════════════════════════════════════════════════════
\section{Configuration dans l'interface AWX}
% ═════════════════════════════════════════════════════════════════════════════

Une fois connecté à AWX, il faut configurer plusieurs ressources dans un ordre précis.

\begin{itemize}
  \item \textbf{Credential Type} : définit la structure d'un type de credential personnalisé (Proxmox n'est pas fourni nativement).
  \item \textbf{Credentials} : les secrets d'authentification (token Proxmox, accès GitHub).
  \item \textbf{Project} : la connexion vers le dépôt Git contenant les playbooks.
  \item \textbf{Inventory} : la liste des cibles sur lesquelles Ansible agit.
  \item \textbf{Job Templates} : l'assemblage final qui relie playbook, inventory et credentials.
\end{itemize}

\begin{infobox}
  \textbf{Note.} Selon la version d'AWX déployée, l'interface peut afficher \textit{Resources \(\to\) Templates} ou \textit{Automation Execution \(\to\) Templates}. La logique reste identique dans les deux cas.
\end{infobox}

\subsection{Création du Custom Credential Type pour Proxmox}

AWX fournit nativement les types courants (Machine SSH, Source Control, etc.) mais pas Proxmox. Définir un type sur mesure permet de stocker proprement les champs spécifiques et de les injecter automatiquement comme variables Ansible dans les playbooks.

\subsubsection*{Création du type}

Naviguer dans \textbf{Administration \(\to\) Credential Types \(\to\) Add} et renseigner :

\begin{itemize}
  \item \textbf{Name} : \cmd{Proxmox API}
  \item \textbf{Description} : \textit{Authentification par token API pour les opérations Proxmox}
\end{itemize}

\subsubsection*{Input configuration}

Cette section définit les champs visibles lors de la création du credential :

\begin{yamlcode}
fields:
  - id: pve_api_host
    type: string
    label: Proxmox host (IP ou FQDN)
  - id: pve_api_user
    type: string
    label: Utilisateur (ex. root@pam)
  - id: pve_api_token_id
    type: string
    label: Token ID (ex. awx)
  - id: pve_api_token_secret
    type: string
    label: Token secret
    secret: true
required:
  - pve_api_host
  - pve_api_user
  - pve_api_token_id
  - pve_api_token_secret
\end{yamlcode}

L'option \cmd{secret: true} garantit que la valeur est masquée dans l'interface et chiffrée en base de données.

\subsubsection*{Injector configuration}

Cette section définit comment les champs sont transmis aux playbooks via \cmd{extra\_vars} :

\begin{yamlcode}
extra_vars:
  pve_api_host:         '{{ pve_api_host }}'
  pve_api_user:         '{{ pve_api_user }}'
  pve_api_token_id:     '{{ pve_api_token_id }}'
  pve_api_token_secret: '{{ pve_api_token_secret }}'
\end{yamlcode}

Ces quatre variables sont ensuite référencées directement dans tous les playbooks et fichiers de tâches. Cliquer sur \textbf{Save}.

\subsection{Création des Credentials}

\subsubsection*{Credential Proxmox}

Naviguer dans \textbf{Resources \(\to\) Credentials \(\to\) Add} :

\begin{itemize}
  \item \textbf{Name} : \cmd{Proxmox AWX Token}
  \item \textbf{Credential Type} : \cmd{Proxmox API}
  \item \textbf{Proxmox host} : \cmd{172.16.1.61}
  \item \textbf{Utilisateur} : \cmd{root@pam}
  \item \textbf{Token ID} : \cmd{awx}
  \item \textbf{Token secret} : la valeur UUID copiée depuis Proxmox
\end{itemize}

Cliquer sur \textbf{Save}.

\subsubsection*{Credential GitHub (dépôt privé)}

Si le dépôt est privé, créer un second credential :

\begin{itemize}
  \item \textbf{Name} : \cmd{GitHub Cloud Playbook}
  \item \textbf{Credential Type} : \cmd{Source Control}
  \item \textbf{Username} : identifiant GitHub
  \item \textbf{Password} : un \textit{Personal Access Token} GitHub généré dans \textit{GitHub \(\to\) Settings \(\to\) Developer settings \(\to\) Personal access tokens \(\to\) Tokens (classic)} avec le scope \cmd{repo}. Ne pas utiliser le mot de passe GitHub, qui n'est plus accepté pour l'authentification API.
\end{itemize}

\begin{infobox}
  \textbf{Note.} Le PAT GitHub est chiffré dans AWX et ne sera plus jamais réaffiché après création, au même titre que le secret Proxmox.
\end{infobox}

\subsection{Création du Project}

Naviguer dans \textbf{Resources \(\to\) Projects \(\to\) Add} :

\begin{itemize}
  \item \textbf{Name} : \cmd{Cloud Playbook}
  \item \textbf{Organization} : \cmd{Default}
  \item \textbf{Source Control Type} : \cmd{Git}
  \item \textbf{Source Control URL} : URL HTTPS du dépôt (ex. \cmd{https://github.com/user/cloud-playbook.git})
  \item \textbf{Source Control Branch} : \cmd{main}
  \item \textbf{Source Control Credential} : \cmd{GitHub Cloud Playbook} (laisser vide si dépôt public)
  \item \textbf{Options} à cocher :
  \begin{itemize}
    \item \textbf{Clean} : nettoie les modifications locales avant chaque synchronisation.
    \item \textbf{Update Revision on Launch} : récupère automatiquement la dernière version du dépôt à chaque lancement de job.
  \end{itemize}
\end{itemize}

Cliquer sur \textbf{Save}. AWX déclenche immédiatement une synchronisation initiale dont le statut doit passer en \textit{Successful} sous une dizaine de secondes.

\subsection{Création de l'Inventory}

L'inventory est la liste des cibles Ansible. Particularité importante : les playbooks de déploiement n'agissent pas par SSH sur Proxmox --- ils appellent l'API REST Proxmox \textit{depuis} le conteneur AWX. La cible Ansible est donc \cmd{localhost} (le conteneur lui-même), même si les commandes atteignent Proxmox via le réseau.

Naviguer dans \textbf{Resources \(\to\) Inventories \(\to\) Add \(\to\) Add inventory} :

\begin{itemize}
  \item \textbf{Name} : \cmd{Infrastructure Proxmox}
  \item \textbf{Organization} : \cmd{Default}
\end{itemize}

Cliquer sur \textbf{Save}, puis dans l'onglet \textbf{Hosts \(\to\) Add} :

\begin{itemize}
  \item \textbf{Name} : \cmd{172.16.1.61}
  \item \textbf{Variables} (YAML) :
\end{itemize}

\begin{yamlcode}
ansible_connection: local
\end{yamlcode}

La directive \cmd{ansible\_connection: local} indique à Ansible de ne pas ouvrir de session SSH : les tâches s'exécutent dans le conteneur AWX, qui appelle ensuite l'API Proxmox via HTTPS.

\begin{infobox}
  \textbf{Note.} Utiliser l'IP \cmd{172.16.1.61} plutôt que \cmd{localhost} améliore la lisibilité des journaux AWX : on identifie clairement que la cible logique est le serveur Proxmox.
\end{infobox}

\subsection{Création des Job Templates}

Il y a quatre Job Templates à créer, un par playbook. Procédure identique pour chacun : \textbf{Resources \(\to\) Templates \(\to\) Add \(\to\) Add job template}.

\subsubsection*{Déployer une VM étudiant (\file{deploy.yml})}

\begin{itemize}
  \item \textbf{Name} : \cmd{Déployer une VM étudiant}
  \item \textbf{Job Type} : \cmd{Run}
  \item \textbf{Inventory} : \cmd{Infrastructure Proxmox}
  \item \textbf{Project} : \cmd{Cloud Playbook}
  \item \textbf{Playbook} : \cmd{deploy.yml}
  \item \textbf{Credentials} : \cmd{Proxmox AWX Token}
\end{itemize}

\subsubsection*{Installer les services applicatifs (\file{services.yml})}

\begin{itemize}
  \item \textbf{Name} : \cmd{Installer les services applicatifs}
  \item \textbf{Playbook} : \cmd{services.yml}
  \item \textbf{Inventory} : \cmd{Infrastructure Proxmox}
  \item \textbf{Credentials} : \cmd{Proxmox AWX Token}
\end{itemize}

\subsubsection*{Démarrer une VM (\file{power.yml})}

\begin{itemize}
  \item \textbf{Name} : \cmd{Démarrer une VM (self-service)}
  \item \textbf{Playbook} : \cmd{power.yml}
  \item \textbf{Inventory} : \cmd{Infrastructure Proxmox}
  \item \textbf{Credentials} : \cmd{Proxmox AWX Token}
\end{itemize}

\subsubsection*{Supprimer une VM (\file{suppression.yml})}

\begin{itemize}
  \item \textbf{Name} : \cmd{Supprimer une VM étudiant}
  \item \textbf{Playbook} : \cmd{suppression.yml}
  \item \textbf{Inventory} : \cmd{Infrastructure Proxmox}
  \item \textbf{Credentials} : \cmd{Proxmox AWX Token}
\end{itemize}

\subsection{Dynamisation via les Surveys}

Les Surveys permettent à chaque étudiant de paramétrer son déploiement via un formulaire. Les réponses sont injectées comme variables Ansible au lancement du job.

Sur la page du Job Template, onglet \textbf{Survey \(\to\) Add}. Créer les trois questions suivantes.

\medskip
\textbf{Question 1 --- Identifiant de la VM}
\begin{itemize}
  \item \textbf{Answer variable name} : \cmd{vmid\_etudiant}
  \item \textbf{Answer type} : \cmd{Integer}
  \item \textbf{Minimum} : \cmd{10} \quad \textbf{Maximum} : \cmd{99}
  \item Obligatoire.
\end{itemize}

\medskip
\textbf{Question 2 --- Système d'exploitation}
\begin{itemize}
  \item \textbf{Answer variable name} : \cmd{choix\_os}
  \item \textbf{Answer type} : \cmd{Multiple Choice (single select)}
  \item \textbf{Choices} : \cmd{UbuntuServ}, \cmd{Debian}, \cmd{Win10}
  \item Obligatoire.
\end{itemize}

\medskip
\textbf{Question 3 --- Services applicatifs}
\begin{itemize}
  \item \textbf{Answer variable name} : \cmd{choix\_services}
  \item \textbf{Answer type} : \cmd{Multiple Choice (multiple select)}
  \item \textbf{Choices} : \cmd{Apache}, \cmd{MariaDB}
  \item Optionnel.
\end{itemize}

Activer le Survey avec le toggle \textbf{Survey Enabled}. Ce même Survey est à reproduire pour \file{services.yml}, \file{power.yml} et \file{suppression.yml} (au minimum les questions \cmd{vmid\_etudiant} et \cmd{choix\_os}).

\subsection{Token d'accès pour le portail web}

L'objectif final est qu'un étudiant passe par un portail web simplifié qui appelle l'API AWX en arrière-plan. Pour cela, générer un token d'accès dans \textbf{User \(\to\) Tokens \(\to\) Add} :

\begin{itemize}
  \item \textbf{Application} : laisser vide pour un token personnel
  \item \textbf{Description} : \cmd{Token portail web}
  \item \textbf{Scope} : \cmd{Write}
\end{itemize}

\textbf{Copier immédiatement} la valeur affichée. L'ID du Job Template est visible dans son URL :\\
\cmd{http://<IP\_AWX>:<PORT>/\#/templates/job\_template/\textbf{7}/details}

Le portail envoie ensuite un \cmd{POST} sur :

\begin{bashcode}
POST http://<IP_AWX>:<PORT>/api/v2/job_templates/7/launch/
Authorization: Bearer <TOKEN_AWX>
Content-Type: application/json

{
  "extra_vars": {
    "vmid_etudiant": 42,
    "choix_os": "UbuntuServ",
    "choix_services": ["Apache", "MariaDB"]
  }
}
\end{bashcode}

AWX renvoie un identifiant de job que le portail peut interroger pour suivre l'avancement via \cmd{/api/v2/jobs/<id>/}.


% ═════════════════════════════════════════════════════════════════════════════
\section{Préparation des templates de VMs}
% ═════════════════════════════════════════════════════════════════════════════

Avant qu'une VM puisse être clonée automatiquement, son image source --- la \textit{template} Proxmox --- doit être correctement préparée. L'objectif est de produire des templates génériques : fonctionnelles immédiatement après clonage, mais suffisamment neutres pour que Cloud-Init (Linux) ou CloudBase-Init (Windows) puissent personnaliser chaque instance.

\subsection{Template Linux (Ubuntu Server / Debian)}

\subsubsection*{Installation de base}

\begin{enumerate}
  \item Créer une VM sur Proxmox avec une ISO Ubuntu Server 22.04 LTS ou Debian 12.
  \item Procéder à une installation minimale : un seul utilisateur \cmd{etu} avec mot de passe \cmd{etu}, partitionnement automatique, OpenSSH activé.
  \item Vérifier que la VM est joignable en SSH.
\end{enumerate}

\subsubsection*{Mise à jour et installation des agents}

\begin{bashcode}
# Mise à jour complète
sudo apt update && sudo apt upgrade -y

# Cloud-Init (généralement pré-installé sur Ubuntu Server)
sudo apt install cloud-init cloud-initramfs-growroot -y

# Agent QEMU pour l'intégration Proxmox
sudo apt install qemu-guest-agent -y
sudo systemctl enable --now qemu-guest-agent
\end{bashcode}

\subsubsection*{Configuration de Cloud-Init}

Cloud-Init doit être configuré pour utiliser uniquement la datasource \cmd{NoCloud} fournie par Proxmox via un disque ConfigDrive :

\begin{bashcode}
# Forcer la datasource NoCloud uniquement
sudo dpkg-reconfigure cloud-init
# Cocher uniquement "NoCloud" dans la liste interactive

# Nettoyage des informations résiduelles
sudo cloud-init clean --logs --seed
\end{bashcode}

\subsubsection*{Configuration de sudo sans mot de passe}

Pour qu'Ansible puisse exécuter des commandes en mode \cmd{become} sans saisie de mot de passe :

\begin{bashcode}
echo "etu ALL=(ALL) NOPASSWD:ALL" | sudo tee /etc/sudoers.d/90-etu
sudo chmod 0440 /etc/sudoers.d/90-etu
\end{bashcode}

\subsubsection*{Préparation finale}

\begin{bashcode}
# Supprimer les clés SSH du serveur (régénérées au premier boot)
sudo rm -f /etc/ssh/ssh_host_*

# Réinitialiser le machine-id
sudo truncate -s 0 /etc/machine-id
sudo rm -f /var/lib/dbus/machine-id
sudo ln -s /etc/machine-id /var/lib/dbus/machine-id

# Vider les journaux et l'historique
sudo cloud-init clean --logs
sudo rm -rf /var/log/*
sudo history -c

# Éteindre proprement
sudo shutdown -h now
\end{bashcode}

\subsubsection*{Conversion en template sur Proxmox}

\begin{enumerate}
  \item Clic droit sur la VM éteinte dans l'interface Proxmox \(\to\) \textbf{Convert to template}.
  \item Ajouter un Cloud-Init Drive : \textbf{Hardware \(\to\) Add \(\to\) CloudInit Drive}, choisir un stockage disponible.
  \item Configurer les valeurs par défaut dans l'onglet \textbf{Cloud-Init} : utilisateur \cmd{etu}, mot de passe \cmd{etu}. L'IP sera injectée dynamiquement par le playbook.
  \item Noter l'ID de la template (affiché dans l'interface) et le reporter dans la map \cmd{templates\_os} de \file{deploy.yml}.
\end{enumerate}

\begin{infobox}
  \textbf{Note.} Les templates Linux utilisent les ID 106 (UbuntuServ) et 103 (Debian) dans le projet actuel. Toute nouvelle template doit être déclarée dans la map \cmd{templates\_os} de \file{deploy.yml}.
\end{infobox}

\subsection{Template Windows (Win10)}

La préparation d'une template Windows est sensiblement plus complexe : elle nécessite l'installation et la configuration manuelle de CloudBase-Init, WinRM, RDP et les pilotes VirtIO.

\subsubsection*{Installation de base}

\begin{enumerate}
  \item Créer une VM Proxmox avec une ISO Windows 10. Ajouter également l'ISO \cmd{virtio-win.iso} comme second lecteur optique pour les pilotes.
  \item Durant l'installation Windows, charger les pilotes VirtIO (stockage et réseau) pour des performances optimales.
  \item Créer le compte utilisateur \cmd{etu} avec mot de passe \cmd{etu}.
  \item Désactiver les mises à jour Windows automatiques.
\end{enumerate}

\subsubsection*{Activation de RDP}

\begin{bashcode}
# === PowerShell en mode Administrateur ===

# Activer RDP dans le registre
Set-ItemProperty `
  -Path 'HKLM:\System\CurrentControlSet\Control\Terminal Server' `
  -Name "fDenyTSConnections" -Value 0

# Démarrer le service en mode automatique
Set-Service -Name TermService -StartupType Automatic
Start-Service TermService

# Ouvrir le pare-feu pour RDP
Enable-NetFirewallRule -DisplayGroup "Bureau a distance"
Enable-NetFirewallRule -DisplayGroup "Remote Desktop"

# Désactiver NLA pour la compatibilité avec Guacamole
$ts = Get-WmiObject -class "Win32_TSGeneralSetting" `
  -Namespace root\cimv2\terminalservices `
  -Filter "TerminalName='RDP-tcp'"
$ts.SetUserAuthenticationRequired(0)
\end{bashcode}

\subsubsection*{Configuration de WinRM}

WinRM est le protocole de gestion à distance utilisé par Ansible pour Windows. Il doit être configuré en HTTPS sur le port 5986 :

\begin{bashcode}
# Activation rapide de WinRM
Enable-PSRemoting -Force
winrm quickconfig -force

# Autoriser l'authentification Basic et NTLM
winrm set winrm/config/service/auth '@{Basic="true"}'
winrm set winrm/config/service '@{AllowUnencrypted="true"}'

# Générer un certificat auto-signé
$cert = New-SelfSignedCertificate `
  -DnsName $env:COMPUTERNAME `
  -CertStoreLocation Cert:\LocalMachine\My

# Créer le listener HTTPS sur le port 5986
winrm create winrm/config/Listener?Address=*+Transport=HTTPS `
  "@{Hostname=`"$env:COMPUTERNAME`"; CertificateThumbprint=`"$($cert.Thumbprint)`"}"

# Ouvrir les ports dans le pare-feu
New-NetFirewallRule -DisplayName "WinRM HTTPS" -Direction Inbound `
  -LocalPort 5986 -Protocol TCP -Action Allow
New-NetFirewallRule -DisplayName "WinRM HTTP"  -Direction Inbound `
  -LocalPort 5985 -Protocol TCP -Action Allow

# Vérifier
winrm enumerate winrm/config/listener
Get-Service WinRM
\end{bashcode}

\begin{warnbox}
  \textbf{Attention.} Sans WinRM correctement configuré dans la template, Ansible ne pourra jamais se connecter aux clones Windows, rendant impossible toute installation ultérieure de services. C'est l'étape la plus critique de la préparation Windows.
\end{warnbox}

\subsubsection*{Installation de CloudBase-Init}

CloudBase-Init permet à la VM de récupérer sa configuration (nom d'hôte, mot de passe, IP) depuis le ConfigDrive Proxmox au démarrage.

\begin{enumerate}
  \item Télécharger CloudBase-Init depuis \cmd{https://cloudbase.it/cloudbase-init/}.
  \item Lancer l'installateur en cochant \textit{Use metadata password} et \textit{Run Cloudbase-Init service as LocalSystem}.
  \item À la fin de l'installation, \textbf{ne pas} cocher « Run Sysprep » : cette étape est effectuée manuellement.
\end{enumerate}

Vérifier que le fichier de configuration active les datasources Proxmox-compatibles :

\begin{bashcode}
# C:\Program Files\Cloudbase Solutions\
#   Cloudbase-Init\conf\cloudbase-init.conf
metadata_services=cloudbaseinit.metadata.services.configdrive.ConfigDriveService,
  cloudbaseinit.metadata.services.nocloudservice.NoCloudConfigDriveService
\end{bashcode}

\subsubsection*{Sysprep et conversion en template}

Sysprep généralise l'installation Windows en supprimant les identifiants uniques (SID, nom d'hôte) de sorte que chaque clone obtienne une identité distincte.

\begin{bashcode}
# Désactiver le verrouillage de compte
net accounts /lockoutthreshold:0

# Lancer Sysprep
C:\Windows\System32\Sysprep\sysprep.exe /generalize /shutdown /oobe `
  /unattend:"C:\Program Files\Cloudbase Solutions\Cloudbase-Init\conf\Unattend.xml"
\end{bashcode}

Une fois la VM éteinte par Sysprep, la convertir en template dans Proxmox (clic droit \(\to\) \textbf{Convert to template}). Ajouter un CloudInit Drive dans l'onglet Hardware. L'ID de la template (107 dans le projet actuel) doit être reporté dans \file{deploy.yml}.

\begin{infobox}
  \textbf{Note.} Sans Sysprep, tous les clones Windows partageraient le même SID, causant des conflits d'identité et empêchant CloudBase-Init de personnaliser correctement chaque instance.
\end{infobox}


% ═════════════════════════════════════════════════════════════════════════════
\section{Référence des playbooks}
% ═════════════════════════════════════════════════════════════════════════════

Cette section documente l'ensemble des fichiers du dépôt Git. Les playbooks s'articulent autour de quatre opérations : déploiement d'une VM, installation de services, démarrage et suppression.

\subsection{Organisation des fichiers}

\begin{itemize}
  \item \file{deploy.yml} --- Playbook principal : clone la VM, configure l'IP et crée la connexion Guacamole.
  \item \file{services.yml} --- Installe Apache et/ou MariaDB sur la VM créée.
  \item \file{power.yml} --- Démarre une VM existante (libre-service étudiant).
  \item \file{suppression.yml} --- Arrête, supprime la VM et nettoie Guacamole.
  \item \file{tasks/create\_linux.yml} --- Clonage, injection d'IP et démarrage d'une VM Linux.
  \item \file{tasks/create\_windows.yml} --- Clonage, injection d'IP et démarrage d'une VM Windows.
  \item \file{tasks/guac\_linux.yml} --- Création de la connexion SSH dans Guacamole.
  \item \file{tasks/guac\_windows.yml} --- Création de la connexion RDP dans Guacamole.
  \item \file{collections/requirements.yml} --- Déclaration des collections Ansible requises.
\end{itemize}

\subsection{\texttt{collections/requirements.yml} --- Dépendances Ansible}

Lu automatiquement par AWX avant chaque exécution :

\begin{yamlcode}
---
collections:
  - name: community.proxmox
  - name: scicore.guacamole
  - name: community.docker
  - name: chocolatey.chocolatey
  - name: ansible.windows
\end{yamlcode}

\begin{itemize}
  \item \cmd{community.proxmox} : module \cmd{proxmox\_kvm} pour cloner, démarrer, éteindre et supprimer des VMs via l'API Proxmox.
  \item \cmd{scicore.guacamole} : modules \cmd{guacamole\_connection} et \cmd{guacamole\_user} pour piloter le serveur d'accès web.
  \item \cmd{community.docker} : module \cmd{docker\_container} pour lancer les conteneurs Apache et MariaDB sur Linux.
  \item \cmd{chocolatey.chocolatey} : module \cmd{win\_chocolatey} pour installer des paquets sous Windows.
  \item \cmd{ansible.windows} : module \cmd{win\_powershell} pour exécuter des scripts PowerShell.
\end{itemize}

\subsection{\texttt{deploy.yml} --- Orchestrateur principal}

S'exécute entièrement sur \cmd{localhost} (connexion locale à l'API Proxmox) et coordonne le processus en cinq phases.

\subsubsection*{Variables internes}

\begin{yamlcode}
vars:
  guac_url:   "http://172.16.180.250:8080/guacamole"
  guac_admin: "guacadmin"
  guac_pass:  "guacadmin"

  templates_os:
    UbuntuServ: 106
    Debian:     103
    Win10:      107
\end{yamlcode}

\subsubsection*{Déroulement des tâches}

\begin{enumerate}
  \item \textbf{Installation de \cmd{proxmoxer}.} La bibliothèque Python est installée via \cmd{pip} dans le conteneur AWX si elle est absente.

  \item \textbf{Création de la VM.} Selon \cmd{choix\_os}, inclusion dynamique de \file{tasks/create\_linux.yml} ou \file{tasks/create\_windows.yml}. Le filtre \cmd{flatten} sur \cmd{[choix\_os]} permet au Survey de transmettre aussi bien une chaîne qu'une liste sans erreur de type.

  \item \textbf{Enregistrement en mémoire.} Un \cmd{add\_host} inscrit l'IP \cmd{172.16.205.\{\{ vmid\_etudiant \}\}} dans le groupe \cmd{new\_vms} avec les credentials \cmd{etu/etu}.

  \item \textbf{Configuration Guacamole.} Inclusion de \file{tasks/guac\_linux.yml} ou \file{tasks/guac\_windows.yml} pour créer la connexion SSH ou RDP.

  \item \textbf{Création du compte étudiant.} Création de l'utilisateur Guacamole \cmd{etu\_\{\{ vmid\_etudiant \}\}} (mot de passe \cmd{etu}) avec accès aux connexions \cmd{SSH\_Linux\_\{\{ vmid\_etudiant \}\}} et \cmd{RDP\_Windows\_\{\{ vmid\_etudiant \}\}}.
\end{enumerate}

\subsection{\texttt{tasks/create\_linux.yml} --- Clonage d'une VM Linux}

\subsubsection*{Clonage depuis la template}

\begin{yamlcode}
- name: "Cloner la VM Template"
  community.proxmox.proxmox_kvm:
    api_host:         "{{ pve_api_host }}"
    api_user:         "root@pam"
    api_token_id:     "{{ pve_api_token_id }}"
    api_token_secret: "{{ pve_api_token_secret }}"
    node:             "exbasis"
    vmid:             "{{ template_id }}"
    clone:            "VMtest"
    newid:            "200{{ vmid_etudiant }}"
    name:             "VM-{{ vmid_etudiant }}-LINUX"
    full:             false
    state:            present
    ciuser:           "etu"
    cipassword:       "etu"
    net:              '{"net0":"virtio,bridge=vmbr0,firewall=0"}'
    force:            true
    validate_certs:   false
\end{yamlcode}

Le clonage est \textbf{lié} (\cmd{full: false}) : le disque du clone partage les blocs communs avec la template et n'alloue de l'espace que pour les données nouvelles, économisant l'espace disque et rendant le clonage quasi-instantané.

\subsubsection*{Injection de l'IP et régénération Cloud-Init}

Le module \cmd{proxmox\_kvm} ne permet pas de configurer \cmd{ipconfig0} lors d'un clonage. L'IP est injectée après coup via un appel direct à l'API REST, puis le disque Cloud-Init est régénéré :

\begin{yamlcode}
- name: "Injecter l'IP dans la config du clone"
  ansible.builtin.uri:
    url: "https://{{ pve_api_host }}:8006/api2/json/nodes/exbasis/
          qemu/200{{ vmid_etudiant }}/config"
    method: POST
    headers:
      Authorization: "PVEAPIToken=root@pam!{{ pve_api_token_id }}=
                      {{ pve_api_token_secret }}"
    body_format: form-urlencoded
    body:
      ipconfig0: "ip=172.16.205.{{ vmid_etudiant }}/16,gw=172.16.1.213"
    validate_certs: false
    status_code: 200

- name: "[Linux] Regénérer le disque Cloud-Init"
  ansible.builtin.uri:
    url: "https://{{ pve_api_host }}:8006/api2/json/nodes/exbasis/
          qemu/200{{ vmid_etudiant }}/cloudinit"
    method: PUT
    headers:
      Authorization: "PVEAPIToken=root@pam!{{ pve_api_token_id }}=
                      {{ pve_api_token_secret }}"
    validate_certs: false
    status_code: 200
\end{yamlcode}

\subsubsection*{Démarrage et synchronisation}

\begin{yamlcode}
- name: "[Linux] Démarrer la nouvelle VM"
  community.proxmox.proxmox_kvm:
    vmid:  "200{{ vmid_etudiant }}"
    state: started
    node:  "exbasis"

- name: "Attendre que le port SSH soit disponible"
  ansible.builtin.wait_for:
    host:    "172.16.205.{{ vmid_etudiant }}"
    port:    22
    delay:   15
    timeout: 300
    state:   started

- name: "Pause pour laisser Cloud-Init finaliser"
  ansible.builtin.pause:
    seconds: 30
\end{yamlcode}

La pause de 30 secondes après l'ouverture du port SSH laisse Cloud-Init terminer son initialisation (configuration du nom d'hôte, génération des clés SSH hôte, application des paramètres réseau) avant que \file{services.yml} ne prenne le relais.

\subsection{\texttt{tasks/create\_windows.yml} --- Clonage d'une VM Windows}

Même logique que le fichier Linux, avec trois différences notables.

\begin{itemize}
  \item Le VMID est préfixé par \cmd{300} au lieu de \cmd{200} : étudiant 12 \(\to\) VMID 30012.
  \item Le nom de la source de clonage est \cmd{WIN10-TEMPLATE} au lieu de \cmd{VMtest}.
  \item Le port attendu est le \textbf{3389 (RDP)}, avec un délai initial de 30 secondes et un délai d'attente maximal de 600 secondes : Windows est significativement plus lent qu'Ubuntu Server à démarrer et à exécuter CloudBase-Init.
\end{itemize}

\begin{yamlcode}
- name: "Attendre que le port RDP soit disponible"
  ansible.builtin.wait_for:
    host:    "172.16.205.{{ vmid_etudiant }}"
    port:    3389
    delay:   30
    timeout: 600
    state:   started
\end{yamlcode}

\subsection{\texttt{tasks/guac\_linux.yml} et \texttt{tasks/guac\_windows.yml}}

\subsubsection*{Connexion SSH (Linux)}

\begin{yamlcode}
- name: "[Guacamole] Créer la connexion SSH"
  scicore.guacamole.guacamole_connection:
    base_url:        "{{ guac_url }}"
    auth_username:   "{{ guac_admin }}"
    auth_password:   "{{ guac_pass }}"
    connection_name: "SSH_Linux_{{ vmid_etudiant }}"
    protocol:        "ssh"
    hostname:        "172.16.205.{{ vmid_etudiant }}"
    port:            22
    username:        "etu"
    password:        "etu"
    validate_certs:  false
    state:           present
\end{yamlcode}

\subsubsection*{Connexion RDP (Windows)}

\begin{yamlcode}
- name: "[Guacamole] Créer la connexion RDP"
  scicore.guacamole.guacamole_connection:
    base_url:                "{{ guac_url }}"
    auth_username:           "{{ guac_admin }}"
    auth_password:           "{{ guac_pass }}"
    connection_name:         "RDP_Windows_{{ vmid_etudiant }}"
    protocol:                "rdp"
    hostname:                "172.16.205.{{ vmid_etudiant }}"
    port:                    3389
    username:                "etu"
    password:                "etu"
    rdp_security:            "nla"
    rdp_ignore_server_certs: true
    validate_certs:          false
    state:                   present
\end{yamlcode}

L'option \cmd{rdp\_ignore\_server\_certs: true} est indispensable car les VMs Windows utilisent un certificat auto-signé généré lors du Sysprep. Sans cette option, Guacamole refuserait la connexion RDP.

Les deux fichiers créent ensuite l'utilisateur Guacamole \cmd{etu\_\{\{ vmid\_etudiant \}\}} avec accès aux deux types de connexions, de sorte qu'un étudiant conserve la même identité Guacamole quel que soit l'OS déployé.

\subsection{\texttt{services.yml} --- Installation des services applicatifs}
\label{sec:services}

Ce playbook est structuré en deux \textit{plays} successifs.

\subsubsection*{Play 1 --- Préparation (\texttt{hosts: localhost})}

Reconstruit l'entrée de la VM en mémoire Ansible et attend que le port de connexion soit accessible. La méthode de connexion est déterminée dynamiquement selon l'OS :

\begin{yamlcode}
- name: Reconstruire
  add_host:
    name:                 "172.16.205.{{ vmid_etudiant }}"
    groups:               "new_vms"
    ansible_user:         "etu"
    ansible_password:     "etu"
    ansible_connection:   "{{ 'winrm' if (choix_os in
                          ['Win10', 'WinServer']) else 'ssh' }}"
    ansible_port:         "{{ 5986 if (choix_os in
                          ['Win10', 'WinServer']) else 22 }}"
    ansible_winrm_server_cert_validation: ignore
    ansible_winrm_transport: ntlm
    ansible_ssh_common_args: >-
      -o StrictHostKeyChecking=no -o UserKnownHostsFile=/dev/null
\end{yamlcode}

Les options \cmd{StrictHostKeyChecking=no} et \cmd{UserKnownHostsFile=/dev/null} évitent l'erreur \cmd{REMOTE HOST IDENTIFICATION HAS CHANGED} lorsqu'une VM est recréée avec la même IP qu'une précédente.

\subsubsection*{Play 2 --- Installation (\texttt{hosts: new\_vms})}

\textbf{Bloc Linux.} Les services sont déployés sous forme de conteneurs Docker.

\begin{enumerate}
  \item Attente de la fin de Cloud-Init (avec gestion du code de retour 2, cf. section~\ref{sec:problemes}).
  \item Installation de \cmd{docker.io}, \cmd{python3-pip} et \cmd{python3-docker} via \cmd{apt}.
  \item Démarrage et activation de Docker via \cmd{systemd}.
  \item Lancement conditionnel des conteneurs selon \cmd{choix\_services} :
  \begin{itemize}
    \item \cmd{httpd:latest} (Apache) sur le port 80:80, avec le volume \cmd{/var/www/html:/usr/local/apache2/htdocs}.
    \item \cmd{mariadb:latest} sur le port 3306:3306, avec le volume \cmd{/home/etu/mysql\_data:/var/lib/mysql} et la variable d'environnement \cmd{MARIADB\_ROOT\_PASSWORD=etu}.
  \end{itemize}
  \item Vérification : test HTTP sur \cmd{localhost} pour Apache, sondage du port 3306 pour MariaDB.
\end{enumerate}

\textbf{Bloc Windows.} Les services sont installés nativement via le gestionnaire de paquets Chocolatey.

\begin{enumerate}
  \item Installation de \cmd{apache-httpd} via Chocolatey si Apache est demandé.
  \item Ouverture idempotente du port 80 dans Windows Defender Firewall.
  \item Installation de \cmd{mariadb} via Chocolatey si MariaDB est demandé.
  \item Ouverture idempotente du port 3306 dans Windows Defender Firewall.
\end{enumerate}

\subsection{\texttt{power.yml} --- Démarrage en libre-service}

Playbook minimal permettant à un étudiant de redémarrer sa VM sans accès à Proxmox. Il contient deux tâches \cmd{proxmox\_kvm} avec \cmd{state: started}, conditionnées sur \cmd{choix\_os} :

\begin{yamlcode}
- name: Démarrer la VM Linux
  community.proxmox.proxmox_kvm:
    vmid:  "200{{ vmid_etudiant }}"
    state: started
    node:  "exbasis"
  when:          "'UbuntuServ' in choix_os or 'Debian' in choix_os"
  no_log:        true
  ignore_errors: true

- name: Démarrer la VM Windows
  community.proxmox.proxmox_kvm:
    vmid:  "300{{ vmid_etudiant }}"
    state: started
    node:  "exbasis"
  when:          "'Win10' in choix_os or 'WinServer' in choix_os"
  no_log:        true
  ignore_errors: true
\end{yamlcode}

La directive \cmd{no\_log: true} masque le token Proxmox dans les journaux AWX. La directive \cmd{ignore\_errors: true} empêche le job d'échouer si la VM est déjà démarrée.

\subsection{\texttt{suppression.yml} --- Suppression d'un déploiement}

Assure le nettoyage complet d'un déploiement étudiant en trois étapes, toutes avec \cmd{ignore\_errors: true} :

\begin{enumerate}
  \item \textbf{Arrêt forcé} (\cmd{state: stopped, force: true}) : équivalent à une coupure d'alimentation, nécessaire si la VM ne répond plus à un arrêt propre.
  \item \textbf{Suppression de la VM} (\cmd{state: absent}) : supprime définitivement la VM et son disque. Linux sur VMID \cmd{200X}, Windows sur VMID \cmd{300X}.
  \item \textbf{Nettoyage Guacamole} : suppression de la connexion \cmd{SSH\_Linux\_\{\{ vmid\_etudiant \}\}} ou \cmd{RDP\_Windows\_\{\{ vmid\_etudiant \}\}} via \cmd{guacamole\_connection} avec \cmd{state: absent}.
\end{enumerate}

\begin{warnbox}
  \textbf{Attention --- Incohérence de conditions.} Les conditions dans \file{suppression.yml} utilisent les prédicats génériques \cmd{'Windows' in choix\_os} et \cmd{'Linux' in choix\_os}, alors que les autres playbooks comparent les valeurs exactes du Survey (\cmd{'Win10'}, \cmd{'UbuntuServ'}, \cmd{'Debian'}). En particulier, la valeur \cmd{Debian} ne contient pas la sous-chaîne « Linux » : le playbook de suppression ne fonctionnera pas correctement pour les VMs Debian dans l'état actuel. Il convient d'harmoniser ces conditions.
\end{warnbox}


% ═════════════════════════════════════════════════════════════════════════════
\section{Portail web étudiant}
% ═════════════════════════════════════════════════════════════════════════════

Le portail web est l'interface finale du projet : il masque complètement AWX et Proxmox aux étudiants, qui n'ont accès qu'à un tableau de bord simplifié depuis leur navigateur. Il est développé en deux parties indépendantes : un backend Python (FastAPI) et un frontend React/TypeScript, servis ensemble derrière un proxy Nginx.

\begin{infobox}
  \textbf{Note --- État du développement.} Le portail est fonctionnel pour les opérations principales (création, démarrage, suppression de VM, accès console Guacamole). Certaines fonctionnalités restent à compléter, notamment l'administration des utilisateurs depuis l'interface.
\end{infobox}

% ─────────────────────────────────────────────────────────────────────────────
\subsection{Architecture technique}
% ─────────────────────────────────────────────────────────────────────────────

\begin{itemize}
  \item \textbf{Backend} : Python 3, \textbf{FastAPI} (asynchrone), base de données SQLite via SQLAlchemy async. Tourne en tant que service systemd \cmd{portail-cloud}, écouté sur \cmd{127.0.0.1:8000} via Uvicorn.
  \item \textbf{Frontend} : \textbf{React} + TypeScript + Tailwind CSS, compilé par Vite en fichiers statiques.
  \item \textbf{Proxy} : \textbf{Nginx} sert le frontend sur le port 80 et redirige les requêtes \cmd{/api/*} vers le backend.
  \item \textbf{Authentification} : double mode --- SSO via \textbf{Microsoft Entra ID} (OAuth2) ou connexion email/mot de passe pour les environnements sans Azure.
\end{itemize}

Le flux d'une requête est le suivant : navigateur \(\to\) Nginx (port 80) \(\to\) soit les fichiers statiques React, soit le backend FastAPI (port 8000) \(\to\) AWX ou Guacamole.

% ─────────────────────────────────────────────────────────────────────────────
\subsection{Installation automatisée}
% ─────────────────────────────────────────────────────────────────────────────

L'installation s'effectue depuis la VM qui hébergera le portail (une Debian/Ubuntu avec accès réseau vers AWX et Guacamole) en exécutant le script fourni à la racine du dépôt :

\begin{bashcode}
git clone https://github.com/user/PortailCloudEtu.git
cd PortailCloudEtu
chmod +x install.sh
./install.sh
\end{bashcode}

Le script exécute huit étapes successives :

\begin{enumerate}
  \item Mise à jour des paquets système (\cmd{apt-get update}).
  \item Installation des dépendances : Python 3, Node.js 20, Git, Nginx.
  \item Création de l'environnement virtuel Python et installation des dépendances via \cmd{pip install -r requirements.txt}.
  \item Génération du fichier \file{.env} avec une clé secrète auto-générée (\cmd{openssl rand -hex 32}) et l'IP locale détectée automatiquement.
  \item Compilation du frontend (\cmd{npm install \&\& npm run build}).
  \item Création et démarrage du service systemd \cmd{portail-cloud} (exécuté sous l'utilisateur \cmd{www-data}).
  \item Configuration de Nginx : frontend en fichiers statiques, proxy \cmd{/api/} vers le backend.
  \item Vérification finale via un appel à \cmd{/api/sante} (doit retourner HTTP 200).
\end{enumerate}

À la fin de l'installation, l'application est accessible sur \cmd{http://<IP\_DE\_LA\_VM>} avec les identifiants par défaut \cmd{admin@local} / \cmd{admin}.

\begin{warnbox}
  \textbf{Attention.} Le compte \cmd{admin@local / admin} est créé automatiquement au premier démarrage. Changer ce mot de passe immédiatement en production depuis l'interface d'administration.
\end{warnbox}

% ─────────────────────────────────────────────────────────────────────────────
\subsection{Configuration --- fichier \texttt{.env}}
% ─────────────────────────────────────────────────────────────────────────────

Le fichier \file{cloud/backend/.env} centralise toute la configuration. Il est généré automatiquement par \file{install.sh} avec des valeurs par défaut à compléter :

\begin{bashcode}
# Application
NOM_APP="Portail Cloud Etudiant"
ENVIRONNEMENT=production
CLE_SECRETE_SESSION=<genere-par-openssl-rand-hex-32>
URL_BDD=sqlite+aiosqlite:///./portail.db

# Microsoft Entra ID (SSO)
ENTRA_TENANT_ID=votre-tenant-id
ENTRA_CLIENT_ID=votre-client-id
ENTRA_CLIENT_SECRET=votre-client-secret
ENTRA_REDIRECT_URI=http://<IP_PORTAIL>/api/auth/callback
GROUPE_ADMIN=cloud-portail-admins
GROUPE_USER=cloud-portail-etudiants

# AWX
AWX_URL=http://<IP_AWX>:<PORT>
AWX_TOKEN=<token-personnel-awx>
AWX_TEMPLATE_DEPLOY=10
AWX_TEMPLATE_DESTROY=11
AWX_TEMPLATE_POWER=12
AWX_TEMPLATE_SERVICES=13

# Guacamole
GUACAMOLE_URL=http://172.16.180.250:8080/guacamole
GUACAMOLE_USER=guacadmin
GUACAMOLE_PASSWORD=guacadmin

# Quotas et CORS
QUOTA_VM_ETUDIANT=2
ORIGINE_AUT=["http://<IP_PORTAIL>"]
\end{bashcode}

Les quatre valeurs \cmd{AWX\_TEMPLATE\_*} correspondent aux ID des Job Templates créés dans AWX (cf. section~4.5). Pour les retrouver : URL du Job Template dans AWX \(\to\) le nombre dans \cmd{.../job\_template/\textbf{10}/...} est l'ID.

Après toute modification du \file{.env}, relancer le backend :

\begin{bashcode}
sudo systemctl restart portail-cloud
\end{bashcode}

% ─────────────────────────────────────────────────────────────────────────────
\subsection{Modèle de données}
% ─────────────────────────────────────────────────────────────────────────────

La base SQLite contient trois tables.

\subsubsection*{Table \texttt{utilisateurs}}

\begin{itemize}
  \item \cmd{entra\_oid} : identifiant unique Microsoft Entra (null pour les comptes locaux).
  \item \cmd{email} : unique, sert d'identifiant de connexion locale.
  \item \cmd{role} : \cmd{etu} ou \cmd{admin}. Un admin voit et gère toutes les VMs.
  \item \cmd{quota\_vms} : nombre maximum de VMs simultanées autorisées (défaut : 2).
\end{itemize}

\subsubsection*{Table \texttt{machines\_virtuelles}}

\begin{itemize}
  \item \cmd{vmid\_etudiant} : identifiant entre 10 et 199, alloué automatiquement (premier libre). Correspond directement au dernier octet de l'IP (\cmd{172.16.205.X}).
  \item \cmd{systeme\_exploitation} : valeur transmise à AWX (\cmd{UbuntuServ}, \cmd{Debian}, \cmd{Win10}).
  \item \cmd{services\_installes} : liste séparée par des virgules (\cmd{Apache,MariaDB}).
  \item \cmd{statut} : cycle de vie --- \cmd{en\_creation} \(\to\) \cmd{demarree} \(\to\) \cmd{arretee} \(\to\) \cmd{en\_suppression} \(\to\) \cmd{supprimee}. En cas d'échec AWX : \cmd{en\_erreur}.
\end{itemize}

\subsubsection*{Table \texttt{jobs}}

Trace chaque opération AWX lancée par le portail. \cmd{awx\_job\_id} permet de re-synchroniser le statut en interrogeant \cmd{/api/v2/jobs/<id>/} côté AWX. Les types de jobs sont : \cmd{creation}, \cmd{suppression}, \cmd{demarrage}, \cmd{arret}, \cmd{installation\_services}.

% ─────────────────────────────────────────────────────────────────────────────
\subsection{Endpoints de l'API REST}
% ─────────────────────────────────────────────────────────────────────────────

La documentation interactive est disponible sur \cmd{http://<IP\_PORTAIL>/api/docs} (Swagger UI généré automatiquement par FastAPI).

\subsubsection*{Authentification --- \texttt{/api/auth}}

\begin{itemize}
  \item \cmd{GET /api/auth/login} \quad Redirige vers la page de connexion Microsoft Entra.
  \item \cmd{GET /api/auth/callback} \quad Callback OAuth2 : crée la session et redirige vers le frontend.
  \item \cmd{POST /api/auth/login} \quad Connexion par email + mot de passe (alternative au SSO).
  \item \cmd{GET /api/auth/me} \quad Retourne les informations de l'utilisateur connecté.
  \item \cmd{POST /api/auth/logout} \quad Vide le cookie de session.
\end{itemize}

\subsubsection*{Gestion des VMs --- \texttt{/api/vms}}

\begin{itemize}
  \item \cmd{GET /api/vms} \quad Liste les VMs actives de l'utilisateur connecté.
  \item \cmd{POST /api/vms} \quad Crée une VM : vérifie le quota, alloue un vmid, lance le Job Template AWX \cmd{deploy} avec les variables \cmd{vmid\_etudiant}, \cmd{choix\_os}, \cmd{choix\_services} et \cmd{student\_name}.
  \item \cmd{GET /api/vms/\{id\}} \quad Détails d'une VM (l'étudiant ne voit que les siennes).
  \item \cmd{POST /api/vms/\{id\}/demarrer} \quad Lance le Job Template AWX \cmd{power}.
  \item \cmd{POST /api/vms/\{id\}/arreter} \quad Lance le Job Template AWX \cmd{power} avec \cmd{action: stop}.
  \item \cmd{POST /api/vms/\{id\}/annuler} \quad Annule le job AWX en cours via \cmd{/api/v2/jobs/<id>/cancel/}.
  \item \cmd{DELETE /api/vms/\{id\}} \quad Lance le Job Template AWX \cmd{destroy}.
  \item \cmd{GET /api/vms/\{id\}/jobs} \quad Historique des jobs d'une VM avec synchronisation des statuts.
  \item \cmd{GET /api/vms/\{id\}/console} \quad Retourne l'URL Guacamole directe pour ouvrir la console dans le navigateur.
\end{itemize}

% ─────────────────────────────────────────────────────────────────────────────
\subsection{Intégration avec AWX}
% ─────────────────────────────────────────────────────────────────────────────

Le portail communique avec AWX via la classe \cmd{ClientAWX} (\file{services/awx\_clients.py}), qui encapsule trois opérations HTTP sur l'API AWX v2 :

\begin{itemize}
  \item \textbf{Lancer un job} : \cmd{POST /api/v2/job\_templates/<id>/launch/} avec \cmd{extra\_vars} en JSON. Retourne l'ID du job créé.
  \item \textbf{Consulter le statut} : \cmd{GET /api/v2/jobs/<id>/}. Mappe les statuts AWX (\cmd{running}, \cmd{successful}, \cmd{failed}...) vers les statuts internes du portail.
  \item \textbf{Annuler} : \cmd{POST /api/v2/jobs/<id>/cancel/}.
\end{itemize}

La synchronisation des statuts est déclenchée à la demande lors de la consultation de l'historique des jobs (\cmd{GET /api/vms/\{id\}/jobs}). Elle met à jour le statut de la VM en base : par exemple, lorsqu'un job de création passe en \cmd{successful}, la VM passe en \cmd{demarree} et son IP (\cmd{172.16.205.<vmid>}) est enregistrée.

\begin{infobox}
  \textbf{Note --- Mode développement.} Lorsque \cmd{ENVIRONNEMENT=developpement} est défini dans le \file{.env}, le \cmd{ClientAWX} est remplacé par un mock qui simule toujours un succès. Cela permet de développer le frontend sans AWX disponible.
\end{infobox}

% ─────────────────────────────────────────────────────────────────────────────
\subsection{Commandes d'exploitation}
% ─────────────────────────────────────────────────────────────────────────────

\begin{bashcode}
# Suivre les logs en temps réel
sudo journalctl -u portail-cloud -f

# Statut du service
sudo systemctl status portail-cloud

# Redémarrer après modification du .env
sudo systemctl restart portail-cloud

# Réinitialiser la base de données (supprime toutes les données)
./reinitialiser.sh

# Vérifier que l'API répond
curl http://localhost/api/sante
# Réponse attendue : {"statut":"ok","app":"Portail Cloud Etudiant"}
\end{bashcode}


% ═════════════════════════════════════════════════════════════════════════════
\section{Problèmes rencontrés et résolutions}
\label{sec:problemes}
% ═════════════════════════════════════════════════════════════════════════════

Le développement de l'infrastructure a mis en évidence plusieurs difficultés techniques. Leur documentation est utile pour les évolutions futures.

\begin{stepbox}{Problème 1 --- Code de retour non zéro de Cloud-Init}
  La tâche \cmd{cloud-init status --wait} retournait systématiquement le code 2, faisant échouer Ansible. Le diagnostic via \cmd{cloud-init status --long} a révélé un statut \cmd{degraded done} causé par un avertissement de dépréciation (utilisation de \cmd{user:} au lieu de \cmd{users:} dans la configuration). La solution consiste à n'échouer que si la sortie ne contient pas \cmd{status: done} :

\begin{yamlcode}
- name: Attendre la fin de Cloud-Init
  ansible.builtin.raw: cloud-init status --wait
  register: cloudinit_result
  changed_when: false
  failed_when:
    - cloudinit_result.rc != 0
    - "'status: done' not in cloudinit_result.stdout"
\end{yamlcode}
\end{stepbox}

\begin{stepbox}{Problème 2 --- Test MariaDB échouant avec \texttt{mariadb-admin ping}}
  La commande cherchait un socket Unix local (\file{/run/mysqld/mysqld.sock}) absent dans le conteneur Docker, qui n'expose le service qu'en TCP. Solution : remplacer la vérification par un test de port avec \cmd{wait\_for} sur \cmd{127.0.0.1:3306}, plus simple et fiable.
\end{stepbox}

\begin{stepbox}{Problème 3 --- Empreinte SSH obsolète après recréation de VM}
  Lorsqu'une VM est recréée avec la même IP qu'une précédente, SSH refuse la connexion avec le message \cmd{REMOTE HOST IDENTIFICATION HAS CHANGED}. Solution : ajout de \cmd{StrictHostKeyChecking=no} et \cmd{UserKnownHostsFile=/dev/null} dans \cmd{ansible\_ssh\_common\_args}. Acceptable dans un contexte de laboratoire où les VMs sont des ressources temporaires.
\end{stepbox}

\begin{stepbox}{Problème 4 --- Adresse APIPA sur les VMs Windows}
  Au déploiement d'une VM Win10 lorsqu'une VM Linux existait déjà avec la même IP, Windows détectait un conflit ARP et revenait sur l'autoconfiguration APIPA (\cmd{169.254.x.x}). Solution : séparer les plages d'adresses IP par OS, ou imposer une seule VM active par étudiant via une tâche de suppression préalable.
\end{stepbox}

\begin{stepbox}{Problème 5 --- Service Apache non créé sous Windows après installation Chocolatey}
  Le paquet Chocolatey \cmd{apache-httpd} installe les binaires d'Apache mais ne crée pas le service Windows associé. Toute tentative de gestion via \cmd{win\_service} échoue avec \cmd{Service is not installed}. Solution retenue : se contenter de l'installation et de l'ouverture du pare-feu, et documenter la vérification manuelle.
\end{stepbox}

\begin{stepbox}{Problème 6 --- Variable \texttt{MYSQL\_ROOT\_PASSWORD} dépréciée}
  L'image officielle MariaDB a remplacé \cmd{MYSQL\_ROOT\_PASSWORD} par \cmd{MARIADB\_ROOT\_PASSWORD}. L'ancienne variable reste fonctionnelle mais génère un avertissement. Les playbooks utilisent désormais la nouvelle variable.
\end{stepbox}


% ═════════════════════════════════════════════════════════════════════════════
\section{Axes d'amélioration}
% ═════════════════════════════════════════════════════════════════════════════

L'infrastructure couvre le cycle de vie complet d'une VM étudiant. Plusieurs pistes d'amélioration sont identifiées pour les itérations suivantes.

\subsection{Harmonisation des conditions dans \texttt{suppression.yml}}

Remplacer les conditions génériques (\cmd{'Windows' in choix\_os}, \cmd{'Linux' in choix\_os}) par les valeurs exactes du Survey (\cmd{'Win10'}, \cmd{'UbuntuServ'}, \cmd{'Debian'}) pour garantir un comportement cohérent avec tous les OS, notamment Debian.

\subsection{Vérification systématique en fin de déploiement}

Ajouter une tâche de vérification à la fin de \file{deploy.yml} qui confirme que la VM est joignable (SSH ou RDP) et que Cloud-Init s'est terminé correctement, avant de déclarer le job AWX comme réussi.

\subsection{Mécanisme de retour arrière}

En cas d'échec partiel (clonage réussi mais injection d'IP échouée), ajouter un bloc \cmd{rescue} qui supprime proprement la VM créée pour éviter de laisser des ressources orphelines dans Proxmox.

\subsection{Extension des services disponibles}

Ajouter de nouveaux choix au Survey : Nginx, PostgreSQL, Redis, GitLab Runner. Chaque service nécessite une entrée dans le Survey, une tâche Docker dans le bloc Linux et une installation Chocolatey dans le bloc Windows.

\subsection{Portail web}

Le portail web (Python/Flask ou Node.js) appelle l'API AWX pour déclencher les Job Templates. Le token AWX généré en section~4.8 est la seule clé d'accès nécessaire côté portail. AWX expose également \cmd{/api/v2/jobs/<id>/} pour suivre l'avancement d'un job en temps réel.

% ═════════════════════════════════════════════════════════════════════════════
\end{document}
% ═════════════════════════════════════════════════════════════════════════════
